\subsection{The Kallosh pages at the end of the supplied source packet}
P000841 begins with the unambiguous handwritten date \ru{24.9.79} and name \ru{Рената Эрнестовна КАЛЛОШ}: a report on a Cambridge meeting on quantum gravity. It belongs to the academic year 1979--80. P000842, despite its separate handwritten page number \(66\), is its continuation: the numbered equations \((1),(2)\) on P000841 continue as \((3),(4),(5)\), with the same spin labels, zero modes, topological numbers and supergravity discussion. It is neither a continuation of Shabat nor a displaced Belavin page. These two pages are accounted for here as supplied foreign-year evidence, without inventing an additional 1978--79 session or assigning their claims to record460.

The first page records Hawking's space-time-foam program, brief reports attributed to Page/Atiyah and Coleman/Jackiw, four-manifold topology, and Singer/DeWitt on supermanifolds. Its mathematical display begins with
\[
 \Delta[A,B]\Phi[A,B]=0,\qquad S=A+B,\quad T=A-B,
\]
where \(\Phi\) is a multispinor, and continues with the recorded zero-mode difference
\[
 n[S,T]-n[S,T-1]
 =4^{S-1}\left[\left(T-\tfrac12\right)\chi
              +\left(S-\tfrac13\right)\tfrac p2\right].
\]
Capitalization of \(S\) is normalized here only for readability. The source calls \(\chi,p\) Euler and Pontryagin numbers, but does not fix the operator, manifold, characteristic-number normalization or boundary conditions. These are therefore a record of formulas, not a verified universal index theorem.

P000842 records the further expression
\[
 n[A,B]-2n[A\mathbin{\overset{?}{-}}\tfrac12,B]
 =4^{A+B-1}\left(A-\tfrac7{12}\right)
       \left(\chi+\tfrac p2\right),
\]
with the altered sign before \(\tfrac12\) unresolved, followed by
\(m[S,T]=m[S]=4^S m[0]\), \(\Delta\Psi_i=\lambda_i\Psi_i\), and \(B_4=n+m\). The note expressly restricts equation~\((3)\) to a self-dual manifold and discusses cancellation of nonzero modes. The claim about absence of divergences has no recorded loop order or regularization and is not promoted to a general quantum-gravity finiteness theorem. The page then records a proposed supermanifold index theorem, the symbols \(E_7\to SU(8)\) and \(\dim E_7=133\) without a specified real form, and the construction of \(F_4\) from automorphisms of the Hermitian \(3\times3\) octonionic Jordan algebra with product \(\tfrac12(AB+BA)\). Its closing remarks about physicists and a tree-like sketch are contextual material, not further identifiable mathematical propositions. Full fate analysis of this separate dated session belongs to the 1979--80 chapter.
